\section{Inspect one sample}
\begin{frame}{Initial sample: keep the raw counts visible}
\princ{RAW COUNTS + DENOMINATORS}
\begin{center}
\tbl{initial}
\end{center}
\vspace{4pt}
\begin{block}{What the row says}
The sample contains \textbf{\NInitialGames} games, with balanced starts. The win rate is wins divided by games, not a floating impression.
\end{block}
\see{Every percentage on this slide can be rebuilt from the visible wins and games.}
\note{Reveal the table after predictions. Ask: what is A's numerator and denominator? Expected: \NInitialWinsA wins out of \NInitialGames games. Then ask for B; expected: \NInitialWinsB out of \NInitialGames.}
\end{frame}

\begin{frame}[shrink=18]{Read the first result narrowly}
\princ{LIMIT THE CLAIM}
\begin{columns}[T]
\begin{column}{0.50\textwidth}
\begin{block}{Supported by this row}
Under the fixed rules and seed \NInitialSeed, B had more observed wins than A in this sample.
\end{block}
\end{column}
\begin{column}{0.46\textwidth}
\begin{block}{Not supported by this row}
That B is generally better, best against every opponent, or proven by one sample.
\end{block}
\end{column}
\end{columns}
\vspace{6pt}
\promptbox{02}{a}
\note{Say the two blocks as a contrast. Question: which word makes the first statement safe? Expected: ``this sample'' and the fixed assumptions.}
\end{frame}

\begin{frame}{One reason the first result might mislead}
\princ{CHECK THE MODEL}
\begin{itemize}
\item A sample is one finite slice of random game outcomes.
\item The seed makes the slice repeatable; repeatable is not the same as universal.
\item Balanced starts address one fairness concern, not every source of variation.
\item A learned strategy can only use the state and training experience the model gives it.
\end{itemize}
\begin{block}{Useful sentence}
This result is evidence about this bounded experiment, not a proof about all possible games.
\end{block}
\note{Choose one reason for the class to write, then name the others as possible checks. Question: what problem does equal starting position solve, and what problem does it leave? Expected: it helps with starter imbalance; it does not remove sampling variation or model limits.}
\end{frame}

\begin{frame}{Tuesday's limited conclusion}
\princ{SCOPE THE CONCLUSION}
\begin{block}{Exit-trace shape}
In this sample, under these fixed rules, \rule{2.2cm}{0.4pt}, but I am not yet willing to claim \rule{2.2cm}{0.4pt} because \rule{2.2cm}{0.4pt}.
\end{block}
\begin{itemize}
\item Name what was observed.
\item Name the condition that makes it true.
\item Name the missing check or limit.
\end{itemize}
\see{A strong answer can refuse to declare a winner while still making a precise observation.}
\note{Collect one sentence. Question: what must appear after ``because''? Expected: a reason such as one finite sample, seed-to-seed variation, or a model limitation.}
\end{frame}
